% version history / change control
%
\documentclass[5p,final]{elsarticle}
%\documentclass[a4paper]{revtex4-l}
% Preamble 
\usepackage{xurl}
\usepackage{amsfonts}
\usepackage{amssymb}
\usepackage{amsmath}  
\usepackage{amsthm}
\usepackage{enumitem} 
\usepackage{geometry}
\usepackage{graphicx}% Include figure files
\usepackage{algorithm}
\usepackage{algcompatible}
\usepackage{algpseudocode}
\usepackage{xfrac}
\usepackage[colorlinks=true, allcolors=blue]{hyperref}
\usepackage{hyperref}
\usepackage[utf8]{inputenc} 
\usepackage[english]{babel} 
\usepackage{url}
%\usepackage{breakurl}
\usepackage{wrapfig}
\usepackage{graphicx}
\usepackage{subfiles}
\usepackage{xfrac}
\usepackage{multirow}
\usepackage{cleveref}
\usepackage{balance}
\usepackage{svg}
\usepackage{booktabs}
\usepackage{tikz}
\usetikzlibrary{shapes.geometric, arrows}
\usepackage{tikz-3dplot}
\usepackage[squaren]{SIunits}
\usepackage{caption}
%\usepackage{widetext}
%\usepackage{cuted}
%	\captionsetup{labelfont={footnotesize,bf},textfont=footnotesize,format=hang,indention=-0.5cm} %makes nicer indents
%	\captionsetup{justification=raggedright,singlelinecheck=false}
\usepackage[normalem]{ulem}
\usepackage{subcaption}
\usepackage{lipsum}
\usepackage{todonotes}
\usepackage{float}
%\usepackage{subfloat}
%\restylefloat{table} % change tab style if required


\newtheorem{theorem}{Theorem}
\newtheorem{lemma}{Lemma}
\newtheorem{definition}{Definition}

\newcommand{\bc}{\begin{center}}
\newcommand{\ec}{\end{center}}
\newcommand{\be}{\begin{equation}}
\newcommand{\ee}{\end{equation}}
\newcommand{\bea}{\begin{eqnarray}}
\newcommand{\eea}{\end{eqnarray}}
\newcommand{\beq}{\begin{eqnarray*}}
\newcommand{\eeq}{\end{eqnarray*}}
\newcommand{\bv}{\left( \begin{array}{c} }
\newcommand{\ev}{\end{array} \right) }
\newcommand{\F}{{\cal{F}}}
\newcommand{\Om}{\Omega}
\newcommand{\om}{\omega}
\newcommand{\Pp}{ \mathbb{P} }
\newcommand{\R}{\mathbb{R}}
\newcommand{\Z}{\mathbb{Z}}
\newcommand{\E}{\mathbb{E}}
\newcommand{\En}{\mathbb{E}_n}
\newcommand{\Et}{\tilde{\mathbb{E}}}
\newcommand{\Etn}{\tilde{\mathbb{E}}_n}
\newcommand{\borel}{ {\cal{B}}(\mathbb{R}) }
\newcommand{\pspace}{(\Om, \F, \Pp)}
\newcommand{\Pt}{\tilde{\mathbb{P}}}
\newcommand{\unitv}{\hat{\vec e}}
\newcommand{\hata}{\hat {\ve \alpha}}
\newcommand{\mean}{\mathrm{mean}}
\newcommand{\std}{\mathrm{std}}
\newcommand{\Var}{\mbox{Var}}
\newcommand{\Cov}{\mathbb{C}\mathrm{ov}}
\newcommand{\RR}{\varrho _{1,2}}
\newcommand{\argmax}{\mathrm{argmax}}
\newcommand{\argmin}{\mathrm{argmin}}
\newcommand{\N}{\mathbb{N}}
\newcommand{\diag}{\mathrm{diag}}
\renewcommand{\vec}[1]{\mathbf{#1}}
\newcommand{\ve}[1]{\boldsymbol{#1}} % for greek letters
\newcommand{\de}{\partial}
\newcommand{\mcl}[1]{\mathcal{#1}}
\newcommand{\clr}[1]{\textcolor{purple}{#1}}
\newcommand{\clb}[1]{\textcolor{blue}{#1}}
\newcommand{\clg}[1]{\textcolor{green}{#1}}
\newcommand{\clp}[1]{\textcolor{red}{#1}}
\newcommand{\mlm}{\lim_{\substack{\Delta x \rightarrow 0 \\ \Delta t \rightarrow 0}}}
\newcommand{\lmx}{\lim_{\Delta x \rightarrow 0 }}
\newcommand{\lmt}{\lim_{\Delta t \rightarrow 0 }}

%% Citation package
\usepackage{natbib}
\bibliographystyle{elsarticle-harv}
\setlength{\emergencystretch}{2em}
% \setcitestyle{authoryear,open={(},close={)}}

\begin{document}

\title{Correlation emergence and the Epps effect in two coupled limit order books}

\author[unsw-mathstats]{Chris Angstmann}\ead{c.angstmann@unsw.edu.au}
\address[unsw-mathstats]{School of Mathematics and Statistics, University of New South Wales, Sydney, NSW 2052, Australia}
\author[uct-sta]{Tim Gebbie} \ead{tim.gebbie@uct.ac.za}
\address[uct-sta]{Department of Statistical Sciences, University of Cape Town, Rondebosch 7701, Western Cape, South Africa}

\begin{abstract}
We give a unified analytic account of correlation emergence and the Epps effect in two coupled limit order books. The Epps effect is the empirical reduction in measured cross-asset log-return correlation at short aggregation scales, or equivalently the build-up of measured correlation as the aggregation interval increases. The model starts from a fixed-grid discrete random walk for order flow in operational time, with creation, cancellation and diffusion. A pair-trader coupling between the books is introduced at the level of order creation, and calendar time is imposed through separate observation clocks. We clarify how the operational-time model reduces to coupled reaction--diffusion equations with a moving reaction boundary defining the model log-mid-price. Using a regularised local-response representation of the coupling, we derive approximate closed-form expressions for realised correlations as a function of aggregation time. Here the Epps effect is shown to arise from three distinct mechanisms: asynchronous event clocks (subordination), finite coupling response times, and their combination.
\end{abstract}

\begin{keyword}
Limit order book \sep Reaction--diffusion \sep Epps effect \sep Subordination \sep Market microstructure
\end{keyword}

%\keywords{anomalous diffusion; limit order book; Epps effect; non-uniform sampling; multi-agent models}
%\pacs{89.65.Gh, 02.50.Ey, 05.40.Fb}


\maketitle

\section{Introduction} \label{sec:introduction}

Financial markets operate through intricately coupled collections of order books. These can be distinguished as being either lit or latent order books where both can contribute to the dynamics of price emergence. Lit order books contain displayed orders that market participants can observe and trade against. Their transparency supports liquidity by making available orders and quoted prices visible. Lit order books match displayed orders according to rules such as price-time priority \citep{ohara1995}. Hidden or non-displayed orders are orders that have already been submitted to a market, but are not displayed in full. 

The latent order book \citep{toth2011} used here is slightly different and something more fundamental. It represents potential supply and demand from trading intentions that have not yet been expressed as orders but are latent in the broader market. \citet{donier2015} introduced the reaction--diffusion representation of the latent order book. \citet{Gant2022b} implemented a stochastic finite-difference solution \citep{angstmann2019advection,angstmann2016} and calibrated the resulting market model. \citet{DianaGebbie2025} then developed the single-book DTRW construction with anomalous diffusion, non-uniform event times and simulated price impact. \citet{BauerDianaGebbie2024} then coupled two simulated lit order books through a pair-trader mechanism and showed numerically that correlation and an Epps effect emerge from the coupled asynchronous dynamics. 

The Epps effect is the empirical decline in measured cross-asset return correlation as the sampling or aggregation interval is shortened, first documented by \citet{epps1979}. The Epps effect has been observed in various empirical studies of financial markets \citep{toth2009,mastromatteo2011}. Here we give an analytical reaction-front treatment of the clock and coupling-response mechanisms, together with their combination to demonstrate the Epps effect. This is done using a clock separation construction where we first construct the reaction boundary in operational time and then subordinate it to some calendar time.

Correlation can then be shown to emerge from the order-flow coupling between the books. The coupling is phenomenological and represents trader interactions; it is not introduced as a prespecified dependence between observed event types or price processes. In a Hawkes model of order-book events, scalar log-return correlation appears after events are mapped into price changes, while the cross-dependence is specified through the event taxonomy and selected kernels \citep{bacry2015,chang2025discrete}. Our construction places the cross-book mechanism directly with a mesoscopic pair-trader coupling whose price effect passes through the order-density fields and moving reaction fronts.

Because the analytical treatment separates the reaction-boundary dynamics in operational time from the asynchronous clocks through which prices are observed, this identifies the attenuation due to clock overlap separately from the finite response of the pair-trader coupling. This naturally formulates the limiting model by subordination rather than by numerical interpolation\footnote{The non-uniform interpolation scheme used in the earlier simulations \citep{BauerDianaGebbie2024} is linked to the limiting construction in \ref{app:legacy-interpolation}}.

\section{The coupled limit order books} \label{sec:2}

We separate the order-book dynamics from the clocks through which they are observed. Let
\begin{equation}
    x_i=i\Delta x,
    \qquad
    u_n=n\Delta u,
    \qquad n\in\mathbb N,
\end{equation}
where $\Delta x>0$ is the fixed log-price-grid spacing, $\Delta u>0$ is the fixed operational-time step, $u$ is operational time and $t$ is reserved for calendar time. At each operational-time step, orders may be created or cancelled, and existing orders may shift by one log-price level. For a single book, the microscopic dynamics is a discrete-time random walk in log-price space with killing and source terms \citep{angstmann2016,angstmann2019advection,DianaGebbie2025}.

Let $\varphi_n^{i(j)}$ denote the signed order density at grid point $x_i$ for asset $j$ at operational-time index $n$, and set
\begin{equation}
    0\leq r\leq1,
    \qquad
    |F_n^{(j)}|\leq r.
\end{equation}
For the continuum construction take $r>0$. The incoming nearest-neighbour weights are $\tfrac12(r+F_n^{(j)})$ and $\tfrac12(r-F_n^{(j)})$, where $F_n^{(j)}$ is the signed stochastic jump bias. The quantity $K_{n-m}^{(\alpha_u)}$ is the Sibuya memory kernel of order $0<\alpha_u\leq1$ \cite{angstmann2016}, and $\nu_j$ is the cancellation rate in book $j$. The exponent $\alpha_u$ belongs to the order-density dynamics in operational time; the calendar-time clock exponent $\alpha_t$ is introduced separately below.

The total source increment is then:
\begin{equation}
    c_n^{i(j)}
    =s_n^{i(j)}+\delta_n^{i(j)}
    +\sum_{k\ne j}\ell_{\varepsilon,n}^{i(j,k)},
    \label{eq:discrete-total-source}
\end{equation}
where $s^{(j)}$ is the one-book order source, $\delta^{(j)}$ is an external shock and $\ell_\varepsilon^{(j,k)}$ is the ordered-pair source acting on book $j$ in response to book $k$, with regularisation parameter $\varepsilon$. On the fixed operational-time and log-price lattice
\begin{equation}
\begin{aligned}
\varphi^{i(j)}_n
&= \sum_{m=0}^{n-1} K_{n-m}^{(\alpha_u)}
 e^{-\nu_j (u_{n-1}-u_m)} \bigg[
\tfrac{1}{2}\big(r + F^{(j)}_{n-1}\big)\varphi^{i-1(j)}_{m} \\
&\quad + \tfrac{1}{2}\big(r - F^{(j)}_{n-1}\big)\varphi^{i+1(j)}_{m}
- r \varphi^{i(j)}_{m}
\bigg] \\
&\quad + e^{-\nu_j\Delta u}\varphi^{i(j)}_{n-1}
+ c_{n-1}^{i(j)}\Delta u.
\end{aligned}
\label{eq:operational-DTRW-update}
\end{equation}
The common-cancellation-rate case is recovered by setting $\nu_1=\nu_2=\nu$. A non-zero jump bias surviving the continuum scaling would generate an advection term. The density equation used below takes the centred limit $F_n^{(j)}=0$.

\subsection{Coupled source terms} \label{ssec:coupledPDE}

An ordered-pair term $\ell_\varepsilon^{(j,k)}$ changes book $j$ in response to book $k$; single-book quantities retain only the index $j$. Boundary conditions ensuring conservation of the integrated signed density make it convenient to consider lit order-book sources \citep{Gant2022b,DianaGebbie2025,BauerDianaGebbie2024}:
\begin{equation}\label{eq:source-term}
s^{(j)}(x,u) = - \lambda_j \mu_j (x - p^{(j)}(u)) e^{-\mu_j\left(x - p^{(j)}(u)\right)^2}.
\end{equation}
Here $p^{(j)}(u)$ is the reaction-boundary log-price defined below, $\lambda_j$ is the source-strength parameter, and $\mu_j$ is the source-shape parameter. The coupling represents an external pairs trader who observes the two books and places side-selective orders in response to their directed log-mid-price difference. We focus on two books, although the ordered-pair notation extends directly to more assets.

\subsection{Moving reaction boundary and calendar clock}

The operational-time model log-mid-price for book $j$ is the unique reaction-boundary zero crossing
\begin{equation}
    p^{(j)}(u)=\{x:\varphi^{(j)}(x,u)=0\}.
    \label{eq:operational-price}
\end{equation}
We assume constant annihilation rates,
\begin{equation}
    a^{(j)}(x,u)=\nu_j,
\end{equation}
and write the total source as
\begin{equation}
    c^{(j)}(x,u)=s^{(j)}(x,u)+\delta^{(j)}(x,u)
    +\sum_{k\ne j}\ell_\varepsilon^{(j,k)}(x,u).
\end{equation}
Order-book shocks $\delta^{(j)}$ can be used to estimate price impact \citep{DianaGebbie2025} by introducing a shock of size $Q$ and measuring the resulting log-price change. Both lit and latent order-book sources contribute to price emergence: lit sources have vanishing boundary conditions, whereas latent sources may have finite boundary values.

Let $T_j(u)$ be a non-decreasing process mapping operational time into calendar time, and define its inverse clock by
\begin{equation}
    E_j(t)=\inf\{u\geq0:T_j(u)>t\}.
    \label{eq:inverse-clock}
\end{equation}
The calendar-time order density and log-price are defined pathwise by
\begin{equation}
    \Phi^{(j)}(x,t)=\varphi^{(j)}\!\left(x,E_j(t)\right),
    \qquad
    P^{(j)}(t)=p^{(j)}\!\left(E_j(t)\right).
    \label{eq:field-price-subordination}
\end{equation}
Under uniqueness of the zero crossing, $P^{(j)}(t)=\{x:\Phi^{(j)}(x,t)=0\}$. Thus $p^{(j)}(u)$ and $P^{(j)}(t)$ are log-prices; exponentiating them gives the corresponding model price levels. The common-clock case is $E_1=E_2$. Asset-specific clocks allow $E_1\ne E_2$ and are required for a separate clock-overlap attenuation mechanism.


% ==========================================================
\section{From discrete order flow to a continuum limit}
% ==========================================================
The fixed-grid operational-time limit is taken under
\begin{equation}\label{eq:d}
    D_{\alpha_u}
    =\lim_{\substack{\Delta x\to0\\\Delta u\to0}}
    \frac{r}{2}\frac{\Delta x^2}{(\Delta u)^{\alpha_u}},
\end{equation}
where $D_{\alpha_u}$ is the diffusion constant\footnote{With dimensions $[D_{\alpha_u}]=[\mathrm{log\text{-}price}]^2[\mathrm{operational\ time}]^{-\alpha_u}$}. Following \citet{angstmann2016,angstmann2019advection}, let $\varphi_\Delta^{(j)}(x,u)$ be the lattice-dependent approximating function satisfying $\varphi_\Delta^{(j)}(x_i,u_n)=\varphi_n^{i(j)}$. Define the operational-time cancellation survival function $\theta_j(u)=e^{-\nu_ju}$. The continuum order density then satisfies
\begin{equation}
\begin{aligned}
\partial_u \varphi^{(j)}
&= D_{\alpha_u}\partial_x^2\left[
\theta_j(u)\,{}^{\mathrm{RL}}D_u^{1-\alpha_u}
\left(\frac{\varphi^{(j)}}{\theta_j(u)}\right)
\right]\\
&\quad-\nu_j\varphi^{(j)}+c^{(j)}(x,u).
\end{aligned}
\label{eq:operational-continuum-density}
\end{equation}
The survival weighting is required when $0<\alpha_u<1$ and $\nu_j>0$. For $\alpha_u=1$, the memory operator has order zero and Equation~\eqref{eq:operational-continuum-density} reduces to the ordinary operational-time reaction--diffusion equation. For $\nu_j=0$, it reduces to the unweighted fractional diffusion form.

The reaction-boundary log-price is
\begin{equation}
    p^{(j)}(u)=\{x:\varphi^{(j)}(x,u)=0\}.
\end{equation}
This moving boundary separates buy- and sell-dominated regions of the order book. The one-front construction does not resolve the best bid and ask separately. The zero crossing is the model log-mid-price; exponentiating it gives the corresponding model price level, which at this reduced level is also used to approximate the realised transaction price. Its calendar-time projection is obtained from Equation~\eqref{eq:field-price-subordination}.

When $E_j$ is an inverse-stable clock of order $0<\alpha_t\leq1$, a closed calendar-time fractional equation follows directly for a linear operational-time generator. With state-dependent or nonlinear coupling sources, the subordinated field remains well defined by Equation~\eqref{eq:field-price-subordination}, but a closed fractional PDE requires a separate linearisation, weak-limit or closure assumption. The exponent $\alpha_t$ is therefore not inferred from $\alpha_u$. The identification $\alpha_u=\alpha_t$ is thus an optional modelling restriction.

\subsection{Regularised analytic coupling notation}
\label{ssec:regularised-analytic-coupling}

The numerical simulations of \citet{BauerDianaGebbie2024} use a finite-grid, side-selective pair-trader source\footnote{Appendix~\ref{app:first-moment-correspondence} reproduces the lattice equation given in \citet{BauerDianaGebbie2024} and notes the corresponding Julia implementation.}. For the analytical derivation we use a bounded regularised representation of the same directed mechanism. The two sources are not pointwise identical. Their relation is discussed in \ref{app:coupling-equivalence} and \ref{app:first-moment-correspondence}.

Write the local coordinate and directed spread as
\begin{equation}
\begin{aligned}
    y_j(x,u)&:=x-p^{(j)}(u),\\
    z_{jk}(u)&:=p^{(j)}(u)-p^{(k)}(u),
    \qquad z_{kj}(u)=-z_{jk}(u).
\end{aligned}
\end{equation}
The analytic source-shape kernel is written
\begin{equation}
    q_j(y)=-\lambda_j\mu_j ye^{-\mu_j y^2},
    \label{eq:qj-analytic}
\end{equation}
where $\lambda_j>0$ and $\mu_j>0$. The first-moment identity requires $q_j$ to be odd, integrable and to have a finite first moment.

The regularised pair-trader coupling from book $k$ into book $j$ is
\begin{equation}
    \ell_\varepsilon^{(j,k)}(x,u)
    =\gamma_{jk}z_{jk}(u)q_j\!\left(y_j(x,u)\right)
      W\!\left(y_j(x,u),z_{jk}(u);\varepsilon\right),
    \label{eq:regularised-coupling}
\end{equation}
where $\gamma_{jk}\geq0$ is a coupling-strength parameter and $W$ is a bounded side-selection function. A smooth choice \cite{OlssonKreiss2005,XiaoHonmaKono2005} is
\begin{equation}
    W(y,z;\varepsilon)
    =\frac{1}{2}\left[1+\tanh\!\left(\frac{yz}{\varepsilon}\right)\right],
    \qquad \varepsilon>0.
    \label{eq:side-selection-function}
\end{equation}
Since $y$ and $z$ both have units of log-price, $\varepsilon$ has units of log-price squared. For fixed $z\ne0$, the transition width in $y$ is $w_\varepsilon(z)\sim\varepsilon/|z|$. We use the reference-scale parameterisation
\begin{equation}
    \varepsilon=|z|_{\mathrm{ref}}w_{\mathrm{ref}},
    \qquad
    \Delta x\ll w_{\mathrm{ref}}\ll L_{q,j},
    \qquad
    L_{q,j}\sim\mu_j^{-1/2}.
\end{equation}
Thus $z_{jk}>0$ selects the $y_j>0$ side of book $j$, while $z_{jk}<0$ selects the $y_j<0$ side. The corresponding hard-threshold source uses $W_0(y,z)=\mathbf{1}_{\{yz>0\}}$.

The smooth and hard-side continuum sources have the same raw first spatial moment, as shown in \ref{app:first-moment-correspondence}. That appendix states a lattice-matching criterion for the finite-grid source equations; it does not identify that criterion with the Julia implementation. Conversion of the raw moment into a reaction-front response coefficient uses the normalized closure and validity assumptions in \ref{app:coupling-only}.

% ==========================================================
\section{Analytic decomposition of the Epps effect}
% ==========================================================

The operational-time log-prices $p^{(j)}(u)$ and their calendar-time images $P^{(j)}(t)$ are from Equation~\eqref{eq:field-price-subordination}. Now, let $R^{(j)}_{\Delta t}(t)=P^{(j)}(t+\Delta t)-P^{(j)}(t)$ denote the calendar-time log return over aggregation interval $\Delta t$.
Then the realised correlation \cite{BarndorffNielsenShephard2004RealizedCovariation} is
\begin{equation}
 \rho_{\Delta t}=\frac{\Cov(R^{(1)}_{\Delta t},R^{(2)}_{\Delta t})}{
 \sqrt{\Var(R^{(1)}_{\Delta t})\Var(R^{(2)}_{\Delta t})}}. \label{eq:realised-corr}
\end{equation}

Equation \eqref{eq:realised-corr} is the population analogue of the realised-correlation estimator used in the high-frequency realised-covariation literature \citep{BarndorffNielsenShephard2004RealizedCovariation}. When the two assets are observed on non-synchronous event clocks, the relevant overlap structure is closely related to non-synchronous covariance estimation \citep{hayashi2005}.

We write $\rho_u$ for the instantaneous correlation of the operational-time innovations and reserve $\rho_\infty$ for the large-aggregation correlation level of the observed log-price model. These quantities coincide only in the clock-only benchmark after the overlap factor tends to one.

\subsection{Subordination-only mechanism}

If log-prices evolve in operational time but are observed through independent clocks, the correlation is reduced
by limited clock overlap.
One finds
\begin{equation}
 \rho_{\Delta t}^{\mathrm{sub}}=\rho_u \, \mathcal{A}_{12}(\Delta t),
\end{equation}
where $\mathcal{A}_{12}(\Delta t)$ is an overlap factor and $\rho_u$ is the instantaneous correlation of the two Brownian innovations in operational time. For independent Poisson refresh clocks, using the pooled rate in the continuous exponential-lag aggregation form gives the tractable approximation
\begin{equation}
 \rho_{\Delta t}^{\mathrm{sub}}
 \simeq
 \rho_u \left[1-\frac{1-e^{-\lambda_{12}^{\rm clk}\Delta t}}{\lambda_{12}^{\rm clk}\Delta t}\right].
\end{equation}
Since $\mathcal A_{12}(\Delta t)\to1$ as $\Delta t\to\infty$, the clock-only benchmark has $\rho_\infty=\rho_u$.

This is analogous to the continuous exponential-lag aggregation envelope of \citet{toth2009}, but with $\lambda_{12}^{\rm clk}$ interpreted here as the effective clock-overlap rate.

See \ref{app:subordination-only} where, for independent Poisson refresh clocks with rates \(\lambda_1^{\rm clk}\) and \(\lambda_2^{\rm clk}\), the pooled refresh process is Poisson with rate $\lambda_{12}^{\rm clk}=\lambda_1^{\rm clk}+\lambda_2^{\rm clk}$.
Here the waiting time until either book refreshes is the minimum of two independent exponential waiting times and is therefore exponential with the pooled rate. The resulting overlap envelope remains an effective modelling approximation.

\subsection{Coupling-only mechanism}

Under the identity-clock benchmark and the normalized reaction-front projection with the local validity conditions in \ref{app:coupling-only}, the coupling gives the calendar-time log-price approximation
\begin{align}
 dP^{(1)}(t)&=\sigma_1 dB_1(t)-\kappa_{12}z_{12}(t)dt, \\
 dP^{(2)}(t)&=\sigma_2 dB_2(t)+\kappa_{21}z_{12}(t)dt,
 \label{eq:coupling-only}
\end{align}
where $z_{12}(t)=P^{(1)}(t)-P^{(2)}(t)$ and $\kappa_{12},\kappa_{21}>0$. The spread relaxes with rate $\kappa=\kappa_{12}+\kappa_{21}$.
Integrating the resulting cross-covariance kernel gives
\begin{equation}
 \rho_{\Delta t}^{\mathrm{coup}}
 \simeq
 \rho_\infty \left[1-\frac{1-e^{-\kappa\Delta t}}{\kappa\Delta t}\right].
\end{equation}

This again takes on the same envelope form as in \citet{toth2009}. In our derivation, its rate is now the reaction-front spread-relaxation rate $\kappa$.

For the closed two-log-price SDE model, the stationary spread implies unit long-aggregation correlation. In the normalised envelope above, \(\rho_\infty\) is retained as a phenomenological large-scale level for broader log-price dynamics outside that reduced closure. The ordered-pair coefficient \(\kappa_{jk}\) is the effective rate with which coupling-induced order flow from book \(k\) moves the reaction boundary of book \(j\). Since the two boundaries move in opposite directions when the spread is nonzero, both motions reduce the same spread.

\subsection{Combined mechanism}

When both effects are present, a leading-order separability approximation gives
\begin{equation}
 \rho_{\Delta t}^{\mathrm{comb}}
 \simeq
 \rho_\infty \, \mathcal{A}_{12}(\Delta t)
 \left[1-\frac{1-e^{-\kappa\Delta t}}{\kappa\Delta t}\right].
\end{equation}
For fractional clocks the exponential term is replaced by a Mittag--Leffler function \cite{Haubold2011MittagLeffler}. 
In the combined approximation, realised correlation builds up when the books sample overlapping operational time and the coupling-induced reaction-boundary response has had time to transmit the common movement. This product form depends on the separability assumption in \ref{app:combined-subordination-coupling}.

\section{Discussion}
\label{sec:discussion}

We connect three mature, but largely disjoint, parts of the market-microstructure literature: reaction--diffusion order books, reduced-form accounts of the Epps effect, and discrete-event models of trading. Each is well developed in its own right.

Reaction--diffusion models successfully capture liquidity replenishment and price impact in single assets, but typically remain uncoupled across assets and make no claims about cross--correlations \citep{donier2015,toth2011,benzaquen2018}. Conversely, Epps--effect models explain scale--dependent correlation using asynchronous sampling, lead--lag effects, or estimator bias, but do so at the level of observed prices rather than order--flow dynamics \citep{epps1979,lo1990,hayashi2005,mastromatteo2011}.
We connect these strands by embedding correlation emergence directly into a coupled reaction--diffusion representation of multiple order books.

Correlation emerges from explicit order-flow coupling between the books rather than from correlated price noise or a latent multivariate diffusion. A limit-order-book Hawkes model of trade event types does not directly specify a scalar log-return correlation; events must first be mapped into price changes. The cross-dependence is nevertheless structural in the latent event model~\footnote{For example, in the four-dimensional Hawkes price model of \citet{chang2025discrete}, the cross-kernel
$\phi^{(c)}(t) =\alpha^{(c)}e^{-\beta t}\mathbf{1}_{\{t>0\}}$
connects price-move events across the two assets. The co-dependence is not strictly emergent, nor phenomenological.}. Our model places the cross-book mechanism in $\gamma_{jk}$, with its price effect transmitted through the coupled order-density fields and moving reaction fronts and as such is sourced from the phenomenology of trading itself rather than being latent pre-specified relationships between trade or event types.

This is also why Hawkes estimation is conditional on the market implementation: matching-engine and order-handling rules can distort parameter estimates and model identification \citep{JericevichChangGebbie2021}. In addition, structural inference in an open and reflexive financial market remains conditional on any latent model imposed on the data \citep{PolakowGebbieFlint2026} irrespective of the calibration approach. Hawkes estimation remains valid, but an estimated excitation network is not automatically the recovered mechanics of the market. This qualification matters when structural and phenomenological parameters are compared.

The reaction boundary carries the correlation between the books. In the reaction--diffusion model, the log-mid-price is defined implicitly as a moving annihilation boundary between buy and sell densities. By analysing how the coupling perturbs the order density near this boundary, we can show that correlation transmission is filtered by the finite relaxation time of the reaction front. This provides a microstructural explanation for the Epps effect that we think is absent from earlier reduced--form derivations, where correlation attenuation is attributed solely to sampling effects or lead--lag structure \citep{toth2009,chang2021}.
Even perfectly synchronous observation does not eliminate short--horizon correlation decay when the coupling response is slow. 

\citet{toth2009} obtain the same elementary exponential aggregation envelope by aggregating exponentially decaying lagged correlations and relate its characteristic scale to the reaction time of market participants. Here $\lambda_{12}^{\rm clk}$ is an event-clock rate, whereas $\kappa$ is a reaction-front spread-relaxation rate derived from the coupled order-book dynamics. This separates the clock and coupling mechanisms and extends the treatment to the fractional and qualified combined cases.

One caution is that the one-front construction treats the zero crossing as the model log-mid-price. Exponentiating it gives the corresponding model mid-price level and, at this reduced level, an approximation to the transaction price. Its finite width coarse-grains discrete log-price levels and part of the unresolved within-book spread structure, but it does not explicitly resolve best bid and best ask. That extension would require two moving reaction fronts together with nontrivial buy- and sell-side currents. 

Nevertheless, this construction allows us to clarify the relation between discrete event dynamics and their continuum limits. Concretely, the exponent $\alpha_u$ belongs to the operational-time density law, while $\alpha_t$ belongs to the calendar-time observation clock. The restriction $\alpha_u=\alpha_t$ is imposed only when the two waiting-time descriptions are deliberately identified \citep{montroll1965,scalas2006,angstmann2016}.

Separating operational time from the observation clocks distinguishes clock overlap from coupling response and their combined effect. Similar Epps envelopes can therefore have different sources. Earlier work separated asynchrony from lead--lag effects \citep{mastromatteo2011}; our treatment starts from interacting order books and keeps the clock and coupling parameters distinct.

The analysis here also separates finite-grid couplings used in the prior numerical simulations from the regularised coupling used here in the analytical derivation. The numerical model follows existing practice in agent--based \cite{dicksetal2024} and order-book simulations \cite{DianaGebbie2025}, where piecewise rules govern order placement. For analytical tractability, we introduce a local regularisation that preserves the directed mechanism and provides a smooth side-response suitable for linearisation around the symmetric state. 

This regularisation is not already present in the limiting process from the discrete event model to the continuum and is best understood as an approximating construction used to impose the required continuity. This distinction allows the numerical and analytical results to be compared without conflating discretisation artefacts with continuum behaviour. The analysis complements arguments that the Epps effect reflects the discrete-event nature of markets \citep{chang2025discrete} by showing how that discreteness manifests after the diffusion limit and which additional assumptions are required.

% ==========================================================
\section{Conclusion}
% ==========================================================

We obtain an analytical explanation of correlation emergence and the Epps effect from the discrete-to-continuum construction of coupled order books and a regularised representation of the pair-trader coupling. Separating reaction-boundary dynamics in operational time from the asynchronous observation clocks distinguishes attenuation due to incomplete clock overlap from attenuation due to the finite coupling response. This provides an analytical counterpart and explanation to the prior numerical simulation studies \citep{BauerDianaGebbie2024,DianaGebbie2025}.

\section*{Acknowledgements}

We acknowledge conversations with Derick Diana, Byron Jacobs and Dominic Bauer.

\section*{Declaration of generative AI and AI-assisted technologies in the manuscript preparation process}

During the preparation of this work, the author(s) used ChatGPT 5.6 Sol for proof reading, structural checking and notation review. The author(s) reviewed and edited the output as needed and take full responsibility for the content of the published article. 

\bibliography{CoupledOB}

\appendix

\section{Subordination-only derivation}
\label{app:subordination-only}

This appendix derives an Epps-effect curve generated only by non-uniform observation clocks. The operational-time order-density dynamics is defined before the clock is applied. The pair-trader coupling is excluded. The observation clocks introduced here are analytical objects used to isolate calendar-time overlap attenuation. 

\subsection{Field and reaction-boundary time change}

Let $\varphi^{(j)}(x,u)$ solve the operational-time order-density equation and let
\begin{equation}
    p^{(j)}(u)=\{x:\varphi^{(j)}(x,u)=0\}
\end{equation}
be its unique reaction boundary. For an increasing process $T_j(u)$, define the inverse clock
\begin{equation}
    E_j(t)=\inf\{u\geq0:T_j(u)>t\}.
\end{equation}
The calendar-time field is
\begin{equation}
    \Phi^{(j)}(x,t)=\varphi^{(j)}\!\left(x,E_j(t)\right).
\end{equation}
Uniqueness of the zero crossing gives
\begin{equation}
    P^{(j)}(t)=\{x:\Phi^{(j)}(x,t)=0\}=p^{(j)}\!\left(E_j(t)\right).
\end{equation}
This is the subordination used in the body. A closed calendar-time equation for $\Phi^{(j)}$ follows without further closure for a linear operational-time generator. The nonlinear coupled field is defined pathwise by the time change; a closed fractional PDE requires an additional approximation.

\subsection{Assumptions for the log-price reduction}

We use the following assumptions.
\begin{enumerate}
\item[(A1)] A filtered probability space supports the two operational-time Brownian innovations, with instantaneous correlation $\rho_u$.
\item[(A2)] The reaction-boundary log-prices are approximated locally by continuous operational-time semimartingales. Drifts are omitted because their short-window covariance contribution is lower order than the martingale covariance.
\item[(A3)] The inverse observation clocks $E_j$ are nondecreasing, right-continuous and exogenous to the Brownian innovations. Their increments have finite first moments. Stationary increments are assumed only when a stationary expression is used.
\item[(A4)] Conditional on the clocks, the covariance of the two observed increments is determined by the overlap of the sampled operational-time intervals.
\end{enumerate}

\subsection{Operational-time log-price reduction}

For the subordination-only calculation, approximate the two reaction-boundary log-prices in a common operational time by
\begin{align}
    dp^{(1)}(u) &= \sigma_1dB^{(1)}(u),\\
    dp^{(2)}(u) &= \sigma_2dB^{(2)}(u),
\end{align}
with
\begin{equation}
    d\langle B^{(1)},B^{(2)}\rangle_u=\rho_u\,du,
\end{equation}
where $\rho_u$ is the operational-time innovation correlation. For the conditional variance and covariance calculations below, the observation clocks are assumed exogenous to the operational-time Brownian innovations. Calendar-time log-prices are $P^{(j)}(t)=p^{(j)}(E_j(t))$. For $I=(t,t+\Delta t]$, define the sampled operational-time intervals
\begin{equation}
    J_j(I)=\big(E_j(t),E_j(t+\Delta t)\big]
\end{equation}
and their overlap
\begin{equation}
    \Theta_{12}(I)=\left|J_1(I)\cap J_2(I)\right|.
\end{equation}
Also write $\Delta E_j(I)=E_j(t+\Delta t)-E_j(t)$ and $R_{\Delta t}^{(j)}=P^{(j)}(t+\Delta t)-P^{(j)}(t)$. Conditional on the clocks,
\begin{equation}
    R_{\Delta t}^{(j)}
    =\sigma_j\int_{E_j(t)}^{E_j(t+\Delta t)}dB^{(j)}(u).
\end{equation}
Writing these stochastic integrals on the common operational-time axis gives
\begin{align}
    R_{\Delta t}^{(1)}
    &=\sigma_1\int_0^\infty
    \mathbf 1_{J_1(I)}(u)dB^{(1)}(u),\\
    R_{\Delta t}^{(2)}
    &=\sigma_2\int_0^\infty
    \mathbf 1_{J_2(I)}(u)dB^{(2)}(u).
\end{align}
The It\^{o} isometry then gives
\begin{equation}
    \operatorname{Var}(R_{\Delta t}^{(j)}\mid E_j)=\sigma_j^2\Delta E_j(I),
\end{equation}
while
\begin{equation}
    \operatorname{Cov}(R_{\Delta t}^{(1)},R_{\Delta t}^{(2)}\mid E_1,E_2)
    =\rho_u\sigma_1\sigma_2\Theta_{12}(I).
\end{equation}
Taking expectations gives
\begin{equation}
    \rho_{\Delta t}^{\rm sub}=\rho_u\mathcal A_{12}(\Delta t),
\end{equation}
where
\begin{equation}
    \mathcal A_{12}(\Delta t)
    =\frac{\mathbb E[\Theta_{12}(I)]}
    {\sqrt{\mathbb E[\Delta E_1(I)]\mathbb E[\Delta E_2(I)]}}.
\end{equation}
Here $t$ is fixed and suppressed in the notation. Since inverse-clock increments need not be stationary, the overlap factor may also depend on $t$.
If $E_1=E_2$, then $\Theta_{12}(I)=\Delta E_1(I)$ and there is no separate attenuation from incomplete clock overlap.

\subsection{Poisson-refresh approximation}

If the books refresh independently with effective rates $\lambda_1^{\rm clk}$ and $\lambda_2^{\rm clk}$, the waiting time until either book refreshes is exponential with the pooled rate
\begin{equation}
    \lambda_{12}^{\rm clk}=\lambda_1^{\rm clk}+\lambda_2^{\rm clk}.
\end{equation}
The probability that neither book refreshes over a lag $s$ is
\begin{equation}
    \mathbb P(\text{no refresh of either book over lag }s)
    =e^{-\lambda_{12}^{\rm clk}s}.
\end{equation}
Averaging the lag-dependent overlap over the aggregation window gives
\begin{align}
    \mathcal A_{12}^{\rm Pois}(\Delta t)
    &\simeq
    \frac{1}{\Delta t}\int_0^{\Delta t}
    \left(1-e^{-\lambda_{12}^{\rm clk}s}\right)ds\\
    &=1-\frac{1}{\Delta t}\int_0^{\Delta t}
    e^{-\lambda_{12}^{\rm clk}s}ds\\
    &=1-\frac{1-e^{-\lambda_{12}^{\rm clk}\Delta t}}
    {\lambda_{12}^{\rm clk}\Delta t}.
\end{align}
Therefore
\begin{equation}
    \rho_{\Delta t}^{\rm sub}
    \simeq
    \rho_u\left[1-\frac{1-e^{-\lambda_{12}^{\rm clk}\Delta t}}
    {\lambda_{12}^{\rm clk}\Delta t}\right].
\end{equation}
The factor is asymptotically $\lambda_{12}^{\rm clk}\Delta t/2$ as $\Delta t\downarrow0$ and tends to one as $\Delta t\to\infty$. Hence $\rho_\infty=\rho_u$ in this benchmark. The pooled-refresh rate is exact under independent Poisson waiting times; the overlap envelope remains an effective modelling approximation.

\subsection{Fractional clock version}

For an inverse-stable clock of order $0<\alpha_t\leq1$, the exponential survival term is replaced by a Mittag--Leffler survival term. If the operational-time log-price is Brownian, then
\begin{equation}
    \operatorname{Var}[P^{(j)}(t)-P^{(j)}(0)]
    =\sigma_j^2\mathbb E[E_j(t)]
    =\frac{\sigma_j^2t^{\alpha_t}}{\Gamma(1+\alpha_t)},
    \label{eq:fractional-price-variance}
\end{equation}
for the standard inverse-stable normalization. Ordinary calendar-time log-price diffusion is recovered at $\alpha_t=1$. Under this log-price reduction, $0<\alpha_t<1$ gives sublinear variance growth. This result follows from the log-price clock, not from the operational-time density equation alone.

Introduce a fractional clock scale $\lambda_{12,\alpha_t}^{\rm clk}$ with units of inverse time to the power $\alpha_t$ and $\left.\lambda_{12,\alpha_t}^{\rm clk}\right|_{\alpha_t=1}=\lambda_{12}^{\rm clk}$. Then
\begin{equation}
    E_{\alpha_t}(-\lambda s^{\alpha_t})
    =\sum_{n=0}^{\infty}
    \frac{(-\lambda)^ns^{\alpha_t n}}
    {\Gamma(\alpha_t n+1)}.
\end{equation}
Termwise integration gives
\begin{align}
    \frac{1}{\Delta t}\int_0^{\Delta t}
    E_{\alpha_t}(-\lambda s^{\alpha_t})ds
    &=\sum_{n=0}^{\infty}
    \frac{(-\lambda)^n(\Delta t)^{\alpha_t n}}
    {(\alpha_t n+1)\Gamma(\alpha_t n+1)}\\
    &=E_{\alpha_t,2}(-\lambda(\Delta t)^{\alpha_t}).
\end{align}
Hence
\begin{equation}
    \mathcal A_{12}^{(\alpha_t)}(\Delta t)
    \simeq
    1-E_{\alpha_t,2}(-\lambda_{12,\alpha_t}^{\rm clk}(\Delta t)^{\alpha_t}).
\end{equation}
Since $E_{1,2}(-x)=(1-e^{-x})/x$, this reduces to the Poisson-refresh expression at $\alpha_t=1$. The Mittag--Leffler expression is an effective fractional extension of the exponential overlap envelope, not an exact joint-overlap law for arbitrary inverse-stable clocks. The exponent $\alpha_t$ is calibrated from the observation-clock waiting times or calendar-time anomalous scaling. It is not fixed by $\alpha_u$.

\section{Coupling-only derivation}
\label{app:coupling-only}

This appendix derives an approximate Epps curve generated by the finite response time of the regularised pair-trader coupling and states the validity conditions for the reaction-front reduction. The front calculation is carried out in operational time. The coupling-only Epps curve then uses the identity clock $E_j(t)=t$.

\subsection{Assumptions}

We use the following assumptions.
\begin{enumerate}
\item[(B1)] Each signed order-density field is differentiable in operational time and twice differentiable in log-price near its reaction boundary for the deterministic local calculation.
\item[(B2)] Each model log-mid-price is a unique simple zero of its signed order-density field.
\item[(B3)] The pair-trader source is local near the reaction boundary and has a finite first spatial moment.
\item[(B4)] The coupling is weak enough that the front profile and its response mode may be frozen over the short coupling-response calculation.
\item[(B5)] The Brownian terms introduced after the deterministic projection are effective log-price innovations. They are not obtained by applying the ordinary chain rule to the complete stochastic density field.
\end{enumerate}

\subsection{Zero-crossing reduction}

The model log-mid-price of book $j$ is the unique zero crossing
\begin{equation}
    \varphi^{(j)}(p^{(j)}(u),u)=0.
\end{equation}
Assume that the locally frozen density field is differentiable near a simple zero and consider the deterministic finite-variation displacement induced by the coupling. Writing
\begin{equation}
    \mathcal L_j(u)=\partial_x\varphi^{(j)}(p^{(j)}(u),u)\neq0,
\end{equation}
differentiation by the ordinary chain rule gives
\begin{equation}
    \frac{d}{du}p_{\rm coup}^{(j)}(u)
    =-\frac{\left.\partial_u\varphi^{(j)}(p^{(j)}(u),u)\right|_{\rm coup}}
    {\mathcal L_j(u)}.
\end{equation}
This identity is used only for the deterministic coupling-induced displacement of the locally frozen front. It is not asserted as a chain rule for the complete stochastic density field. A stochastic moving boundary would generally require an It\^{o}--Wentzell treatment. The Brownian innovations introduced below are therefore effective log-price terms and are not derived from this identity. The identity depends on the point value of the source. The regularised source is odd locally and satisfies
\begin{equation}
    \ell_\varepsilon^{(j,k)}(p^{(j)}(u),u)=0.
\end{equation}
Its distributed effect on a finite-width reaction front therefore requires a separate collective-coordinate closure; it does not follow from point evaluation alone.

Let $\varphi_j^\star(y)$ be the locally frozen front profile, where $y=x-p^{(j)}(u)$, and define its translational mode by
\begin{equation}
    v_j(y)=-\partial_y\varphi_j^\star(y).
\end{equation}
If $\psi_j$ is the corresponding adjoint response mode, with inner product $\langle\cdot,\cdot\rangle$, set
\begin{equation}
    N_j=\langle\psi_j,v_j\rangle\neq0.
\end{equation}
Projection of the locally linearised density equation onto $\psi_j$ gives the normalized coupling-induced velocity
\begin{equation}
    \frac{d}{du}p_{\rm coup}^{(j)}(u)
    =\frac{\langle\psi_j,\ell_\varepsilon^{(j,k)}\rangle}{N_j}.
    \label{eq:normalised-front-projection}
\end{equation}
The sign is fixed by the definition of $v_j$. If, over the active support of the coupling source,
\begin{equation}
    \frac{\psi_j(y)}{N_j}\simeq A_jy,
    \label{eq:first-moment-mode}
\end{equation}
then Equation~\eqref{eq:normalised-front-projection} reduces to the first-moment closure
\begin{equation}
    \frac{d}{du}p_{\rm coup}^{(j)}(u)
    \simeq A_j\int_{\mathbb R}y\ell_\varepsilon^{(j,k)}(y,z_{jk})\,dy.
    \label{eq:first-moment-front-closure}
\end{equation}
The coefficient $A_j$ contains the normalization, orientation and dimensions required by the projection. If $\varphi^{(j)}$ is measured as signed order volume per unit log-price, then $A_j$ has units of inverse order volume, so Equation~\eqref{eq:first-moment-front-closure} has units of log-price per unit operational time. \ref{app:first-moment-correspondence} proves that the raw moment is $\gamma_{jk}M_j^+z_{jk}$, where
\begin{align}
    M_j^+&=\int_0^\infty yq_j(y)\,dy
    &=-\lambda_j\mu_j\int_0^\infty
    y^2e^{-\mu_jy^2}dy\\
    &=-\frac{\lambda_j\sqrt{\pi}}{4\sqrt{\mu_j}},
\end{align}
using
\begin{equation}
    \int_0^\infty y^2e^{-\mu y^2}dy
    =\frac{\sqrt{\pi}}{4\mu^{3/2}}.
\end{equation}
Thus
\begin{align}
    \frac{d}{du}p_{\rm coup}^{(j)}
    &\simeq A_j\gamma_{jk}M_j^+z_{jk}
    =-\kappa_{jk}z_{jk},\\
    \kappa_{jk}&:=-A_j\gamma_{jk}M_j^+>0.
    \label{eq:ordered-pair-rate}
\end{align}
For the signed-density convention used here we assume $A_j>0$. Since $M_j^+<0$ and $\gamma_{jk}>0$, this gives the positive mean-reverting rate $\kappa_{jk}>0$.

\subsection{Linear response}

Let $z(u)=z_{12}(u)=p^{(1)}(u)-p^{(2)}(u)$. Around $z=0$, the side-selected coupling induces a stabilising linear spread response. The operational-time log-price approximation is
\begin{align}
    dp^{(1)}(u) &= \sigma_1dB_1(u)-\kappa_{12}z(u)du,\\
    dp^{(2)}(u) &= \sigma_2dB_2(u)+\kappa_{21}z(u)du,
\end{align}
where $B_1$ and $B_2$ are independent Brownian motions. Hence
\begin{equation}
    dz(u)=-(\kappa_{12}+\kappa_{21})z(u)du+\sigma_zdB_z(u),
\end{equation}
where
\begin{equation}
    \sigma_z^2=\sigma_1^2+\sigma_2^2,
    \qquad
    \sigma_zdB_z=\sigma_1dB_1-\sigma_2dB_2,
\end{equation}
and $\kappa=\kappa_{12}+\kappa_{21}$ is the operational-time coupling response rate. Thus $B_z$ is constructed from the two log-price innovations rather than introduced as a third independent source. For the coupling-only calculation below, the identity clock identifies $u$ with calendar time $t$.

\subsection{Validity of the local zero-crossing reduction}

For a coupling-induced front displacement $\xi_j$, write
\begin{equation}
    \varphi_j^\star(\xi_j)
    =\mathcal L_j\xi_j
    +\frac12\mathcal C_j\xi_j^2
    +O(\xi_j^3),
    \qquad
    \mathcal C_j=\partial_y^2\varphi_j^\star(0).
\end{equation}
The Taylor remainder is small relative to the retained term when
\begin{equation}
    \eta_j^{\rm disp}
    :=\frac{|\mathcal C_j|\,|\xi_j|}{2|\mathcal L_j|}
    \ll1.
\end{equation}
The simple-root condition $\mathcal L_j\neq0$ must hold, and the coupling-source width $w_{q,j}\sim\mu_j^{-1/2}$ must remain within the locally linear part of the front:
\begin{equation}
    w_{q,j}\ll L_{{\rm curv},j}.
\end{equation}
The frozen-shape approximation also requires the front-shape evolution time to exceed the response time,
\begin{equation}
    \tau_{{\rm shape},j}\gg\kappa^{-1}.
\end{equation}
The reduction may fail for a flat or non-simple zero, multiple zero crossings, large shocks or displacements, strong coupling that deforms rather than translates the front, a source broad relative to the local curvature scale, front evolution on the response timescale, or asymmetric finite-domain and grid effects. Outside this regime the response is more generally
\begin{equation}
    dp_{\rm coup}^{(j)}=b_j(z,u)du,
\end{equation}   
where $b_j(z,u)=-\kappa_{jk}z$ + higher-order + state-dependent terms. The spread is then not generally Ornstein--Uhlenbeck and the exponential response kernel below does not follow directly.

The finite one-front width coarse-grains discrete log-prices and part of the unresolved within-book spread structure. It does not guarantee local linearity or constitute an explicit bid--ask spread. Resolving that spread requires distinct bid and ask fronts with nontrivial buy- and sell-side currents.

\subsection{Integrated covariance kernel}

For the closed two-log-price SDE, define
\begin{equation}
    C
    =\frac{\kappa_{12}^2\sigma_2^2+\kappa_{21}^2\sigma_1^2}
    {(\kappa_{12}+\kappa_{21})^2},
    \qquad
    A_{\Delta t}
    =\Delta t-\frac{1-e^{-\kappa\Delta t}}{\kappa}.
\end{equation}
The finite response time $\kappa^{-1}$ gives the even component of the lagged cross-covariance kernel
\begin{equation}
    c_{12}^{\rm even}(u)=C\frac{\kappa}{2}e^{-\kappa|u|}.
\end{equation}
A parameter-asymmetric model may also contain an odd lead--lag component. It vanishes under the balance condition $\kappa_{21}\sigma_1^2=\kappa_{12}\sigma_2^2$. Because the equal-window aggregation weight is even, the odd component does not contribute to the covariance calculated below.
The covariance of log returns over a window of length $\Delta t$ is
\begin{align}
    \operatorname{Cov}_{\Delta t}
    &=\int_{-\Delta t}^{\Delta t}
    (\Delta t-|u|)c_{12}^{\rm even}(u)du\\
    &=C\kappa\int_0^{\Delta t}
    (\Delta t-u)e^{-\kappa u}du.
\end{align}
The remaining integral is
\begin{align}
    \int_0^{\Delta t}(\Delta t-u)e^{-\kappa u}du
    &=\frac{\Delta t}{\kappa}
    -\frac{1-e^{-\kappa\Delta t}}{\kappa^2},
\end{align}
and therefore
\begin{equation}
    \operatorname{Cov}_{\Delta t}
    =C\left[\Delta t-
    \frac{1-e^{-\kappa\Delta t}}{\kappa}\right]
    =CA_{\Delta t}.
\end{equation}
The exact equal-window log-return variances of the closed SDE are
\begin{equation}
    \operatorname{Var}(R_{\Delta t}^{(j)})
    =\sigma_j^2\Delta t+(C-\sigma_j^2)A_{\Delta t}.
\end{equation}
Its exact correlation is therefore
\begin{equation*}
    \rho_{\Delta t}^{\rm SDE}
    =\frac{CA_{\Delta t}}
    {\sqrt{[\sigma_1^2\Delta t+(C-\sigma_1^2)A_{\Delta t}]
    [\sigma_2^2\Delta t+(C-\sigma_2^2)A_{\Delta t}]}}.
\end{equation*}
At short scales,
\begin{equation}
    \rho_{\Delta t}^{\rm SDE}
    \sim\frac{C\kappa}{2\sigma_1\sigma_2}\Delta t,
    \qquad \Delta t\downarrow0,
\end{equation}
while the stationary spread implies $\rho_{\Delta t}^{\rm SDE}\to1$ as $\Delta t\to\infty$.

For comparison with a broader log-price process containing variance components outside this reduced coupling closure, use the phenomenological approximation $\operatorname{Var}(R_{\Delta t}^{(j)})\simeq v_j\Delta t$. This gives
\begin{equation}
    \rho_{\Delta t}^{\rm coup}
    \simeq
    \rho_\infty
    \left[1-\frac{1-e^{-\kappa\Delta t}}{\kappa\Delta t}\right],
    \qquad
    \rho_\infty=\frac{C}{\sqrt{v_1v_2}}.
\end{equation}
Here $\rho_\infty$ is an effective large-scale normalisation rather than a free parameter of the closed two-log-price SDE. Setting $\rho_\infty=1$ recovers the long-scale implication of that closed model. The phenomenological expression is linear at the origin and converges to $\rho_\infty$ at large aggregation scales.

\section{Combined derivation}
\label{app:combined-subordination-coupling}

This appendix combines the observation-clock and finite-coupling mechanisms. The earlier numerical construction entangles the non-uniform event mesh with the coupling response because the realised increments enter the evolving book states. The analytical product below instead separates calendar-time clock overlap from operational-time coupling response. It is a leading-order mechanism decomposition, not a pathwise transcription of the numerical update. In operational time the coupling produces an exponential response, while the inverse clock maps that response into calendar time. For inverse-stable subordination of order $\alpha_t$, the response becomes Mittag--Leffler, as in the theory of inverse stable subordinators \citep{MeerschaertStraka2013}.

\subsection{Assumptions}

We use the following assumptions.
\begin{enumerate}
\item[(C1)] In operational time, the coupling-induced covariance response has a completely monotone relaxation survival function $h(u)$ with response density $K(u)=-h'(u)$.
\item[(C2)] The observation clocks are exogenous to the effective log-price innovations and independent of the local coupling response to first order.
\item[(C3)] Variance growth is approximately linear over the aggregation scales used in the reduced expression. Otherwise the corresponding variance normalization must be retained explicitly.
\item[(C4)] The product formula is a leading-order separable approximation. It need not hold when the clock changes the local coupling intensity or pair-trader activity is driven by the same events that generate the observations.
\end{enumerate}

\subsection{Clock-overlap combination}

Let $H_\kappa(I)$ denote the normalized coupling response over the aggregation window. The product form requires the leading-order separability condition
\begin{equation}
    \mathbb{E}[\Theta_{12}(I)H_\kappa(I)]
    \simeq
    \mathbb{E}[\Theta_{12}(I)]\,\mathbb{E}[H_\kappa(I)].
    \label{eq:clock-coupling-separability}
\end{equation}
If the clocks mainly attenuate the operational-time overlap while the coupling response remains approximately exponential in calendar time, this gives
\begin{equation}
    \rho_{\Delta t}^{\rm comb}
    \simeq
    \rho_\infty\mathcal{A}_{12}(\Delta t)
    \left[1-\frac{1-e^{-\kappa\Delta t}}{\kappa\Delta t}\right].
\end{equation}
This form is useful for light-tailed non-uniform clocks because the two attenuation mechanisms remain separately identifiable. If the observation clocks are common, $E_1=E_2$, then $\mathcal A_{12}(\Delta t)=1$ and the expression reduces to the coupling-response curve. This does not remove the non-uniform temporal dependence already embedded in a numerical state update using a common random mesh. With distinct observation clocks, $E_1\ne E_2$, the additional overlap factor remains nontrivial. If clock overlap and coupling transmission are materially dependent, the expectation does not factorise and the product form need not hold.


\subsection{Fractional response}

Let $h(u)$ denote the operational-time relaxation survival function. Ordinary relaxation solves
\begin{equation}
    \frac{d}{du}h(u)=-\kappa h(u),
    \qquad h(0)=1,
\end{equation}
with $h(u)=e^{-\kappa u}$. Under inverse-stable subordination of order $0<\alpha_t\leq1$, introduce a calendar-time response scale $\kappa_{\alpha_t}$ with units of inverse time to the power $\alpha_t$ and $\left.\kappa_{\alpha_t}\right|_{\alpha_t=1}=\kappa$. The subordinated scalar relaxation uses the Caputo derivative:
\begin{equation}
    {}^{\mathrm C}D_t^{\alpha_t}h_{\alpha_t}(t)
    =-\kappa_{\alpha_t}h_{\alpha_t}(t),
    \qquad h_{\alpha_t}(0)=1.
\end{equation}
For $0<\alpha_t<1$,
\begin{equation}
    {}^{\mathrm C}D_t^{\alpha_t}h(t)
    =\frac{1}{\Gamma(1-\alpha_t)}
    \int_0^t\frac{h'(s)}{(t-s)^{\alpha_t}}\,ds,
\end{equation}
Taking Laplace transforms of the fractional relaxation equation gives
\begin{equation}
    s^{\alpha_t}\widehat h_{\alpha_t}(s)-s^{\alpha_t-1}
    =-\kappa_{\alpha_t}\widehat h_{\alpha_t}(s),
\end{equation}
and hence
\begin{equation}
    \widehat h_{\alpha_t}(s)
    =\frac{s^{\alpha_t-1}}
    {s^{\alpha_t}+\kappa_{\alpha_t}}.
\end{equation}
Inverting gives
\begin{equation}
    h_{\alpha_t}(t)=E_{\alpha_t}(-\kappa_{\alpha_t}t^{\alpha_t}).
\end{equation}
The corresponding response density is
\begin{align}
    K_{\alpha_t,\kappa_{\alpha_t}}(t)
    &=-\frac{d}{dt}E_{\alpha_t}
    (-\kappa_{\alpha_t}t^{\alpha_t})\\
    &=\kappa_{\alpha_t}t^{\alpha_t-1}
    E_{\alpha_t,\alpha_t}
    (-\kappa_{\alpha_t}t^{\alpha_t}).
\end{align}
Since $K=-h'$, integration by parts gives
\begin{align}
    \int_0^{\Delta t}(\Delta t-s)K(s)ds
    &=\Delta t-\int_0^{\Delta t}h(s)ds.
\end{align}
The normalised covariance build-up factor is
\begin{align}
    F_{\alpha_t}(\Delta t)
    &=\frac{1}{\Delta t}\int_0^{\Delta t}(\Delta t-s)
    K_{\alpha_t,\kappa_{\alpha_t}}(s)ds\\
    &=1-\frac{1}{\Delta t}\int_0^{\Delta t}
    E_{\alpha_t}(-\kappa_{\alpha_t}s^{\alpha_t})ds\\
    &=1-E_{\alpha_t,2}(-\kappa_{\alpha_t}(\Delta t)^{\alpha_t}).
\end{align}
Therefore
\begin{equation}
    \rho_{\Delta t}^{(\alpha_t)}
    \simeq
    \rho_\infty\left[1-E_{\alpha_t,2}
    (-\kappa_{\alpha_t}(\Delta t)^{\alpha_t})\right].
\end{equation}
For $\alpha_t=1$, $E_{1,2}(-x)=(1-e^{-x})/x$, recovering the coupling-only curve. For small $\Delta t$,
\begin{equation}
    E_{\alpha_t,2}(-\kappa_{\alpha_t}(\Delta t)^{\alpha_t})
    =1-\frac{\kappa_{\alpha_t}(\Delta t)^{\alpha_t}}
    {\Gamma(\alpha_t+2)}
    +O((\Delta t)^{2\alpha_t}),
\end{equation}
so
\begin{equation}
    \rho_{\Delta t}^{(\alpha_t)}
    \sim
    \rho_\infty\frac{\kappa_{\alpha_t}(\Delta t)^{\alpha_t}}
    {\Gamma(\alpha_t+2)}.
\end{equation}
The inverse clock changes the initial build-up from linear in $\Delta t$ to order $(\Delta t)^{\alpha_t}$. Identifying $\alpha_t$ with the operational-time density exponent $\alpha_u$ is an additional modelling restriction.

\section{Coupling correspondence}\label{app:coupling-equivalence}

Regularised-interface and immersed-boundary calculations distinguish pointwise equality from weak agreement after projection onto the modes or moments entering a reduced equation \cite{ElderGrantProvatasKosterlitz2001,Peskin2002,OlssonKreiss2005,DesjardinsMoureauPitsch2008}. We use \(W(y,z;\varepsilon)\) in this projected sense. The analytical pair-trader source is not required to coincide pointwise with either the finite-grid implementation or the lattice equation from~\cite{BauerDianaGebbie2024}. What is required for the log-price reduction is a controlled projection onto the reaction-front displacement mode.
The construction is analogous to sharp-interface asymptotics for diffuse-interface models, where a finite-width interface is introduced and the reduced moving-boundary law is obtained by projection before taking the sharp-interface limit~\cite{ElderGrantProvatasKosterlitz2001}.

The comparison below is local and projection-based. 

\subsection{Objects being compared}

Let $p^{(j)}$ and $p^{(k)}$ be the reaction-boundary log-prices of two books, and write
\begin{equation}
y=x-p^{(j)},
\qquad z=z_{jk}=p^{(j)}-p^{(k)}.
\end{equation}
Appendix~\ref{app:first-moment-correspondence} gives both the finite-grid Julia implementation and the piecewise lattice equations from~\citet{BauerDianaGebbie2024}. The analytical construction below is compared through the reduced reaction-front response rather than by matching either source pointwise.

The analytic coupling is
\begin{equation}
\ell_{\rm reg}^{(j,k)}(x,u)
=\gamma_{jk}zq_j(y)W(y,z;\varepsilon),
\label{eq:eq-reg}
\end{equation}
where $q_j$ is odd and locally linear at $y=0$, and $W$ is bounded with hard-side limit
\begin{equation}
W_0(y,z)=\mathbf 1_{\{yz>0\}}.
\end{equation}

\subsection{Pointwise comparison and local matching criterion}

The numerical and analytical couplings use different side-response rules, so pointwise equality of the complete source fields is not expected. For the log-price reduction, the relevant requirement is agreement of the projected first-order response about $z=0$. In the frozen-slope approximation this projection is represented by the first moment below.

Define the projected first moment by
\begin{equation}
\mathcal M[\ell](z)=\int_{\mathbb R}y\ell(y,z)\,dy.
\end{equation}
A finite-grid source and the regularised source share the same leading response at the level relevant for the log-price reduction if
\begin{equation}
\mathcal M[\ell_{\rm disc}^{\Delta x}](z)
=\eta_{jk} z+o(z)
\end{equation}
and
\begin{equation}
\mathcal M[\ell_{\rm reg}^{\varepsilon}](z)
=\eta_{jk} z+o(z)
\label{eq:weak-equiv}
\end{equation}
with the same finite ordered-pair coefficient $\eta_{jk}$, after matching the discrete coupling strength to $\gamma_{jk}$.

For the regularised coupling,
\begin{align}
\mathcal M[\ell_{\rm reg}](z)
&=\gamma_{jk}z\int_{\mathbb R}yq_j(y)W(y,z;\varepsilon)\,dy.
\end{align}
In the hard-side limit, for $z>0$,
\begin{equation}
\mathcal M[\ell_{\rm reg}](z)
=\gamma_{jk}z\int_0^\infty yq_j(y)\,dy
=\gamma_{jk}M_j^+z.
\end{equation}
Thus the regularised coupling satisfies \eqref{eq:weak-equiv} with
\begin{equation}
\eta_{jk}=\gamma_{jk}M_j^+.
\end{equation}
For the Gaussian source shape,
\begin{equation}
M_j^+=-\frac{\lambda_j\sqrt\pi}{4\sqrt{\mu_j}}.
\end{equation}
A finite-grid source has the same leading reduced response if its lattice first moment satisfies
\begin{equation}
\Delta x\sum_i (x_i-p^{(j)})\ell_{{\rm disc},i}^{(j,k)}
=\gamma_{jk}M_j^+z+o(z)+o(1)
\qquad
\label{eq:lattice-moment-match}
\end{equation}
where $\Delta x\to0$ is taken first and $z\to0$ second, or the analogous joint limit with a specified relation between $z$, $\Delta x$, and $\varepsilon$.

\subsection{The regularisation parameter}

The smooth selector
\begin{equation}
W(y,z;\varepsilon)=\frac12\left[1+\tanh\left(\frac{yz}{\varepsilon}\right)\right]
\end{equation}
introduces a transition layer of width approximately $\varepsilon/|z|$ in the $y$ variable.  If $\varepsilon$ is fixed and $z\to0$, then $W\to1/2$ pointwise, and because $yq_j(y)$ is even,
\begin{equation}
\int_{\mathbb R}yq_j(y)W(y,z;\varepsilon)dy
\to \frac12\int_{\mathbb R}yq_j(y)dy
=\int_0^\infty yq_j(y)dy.
\end{equation}
Thus the first-moment coefficient remains finite and equals the hard-side coefficient for symmetric kernels.  If an asymmetric kernel is used, the order of limits $\varepsilon\to0$ and $z\to0$ must be specified explicitly.

A general reference-scale choice is
\begin{equation}
\varepsilon=|z|_{\rm ref}w_{\rm ref},
\qquad
\Delta x\ll w_{\rm ref}\ll L_{q,j},
\end{equation}
where $L_{q,j}$ is the spatial scale on which $q_j$ varies. The first inequality makes the smoothed side selector resolvable on the lattice, while the second keeps the selector local relative to the order-book source shape. The earlier square-root expression is recovered only under the symmetric reference-scale convention in which $|z|_{\rm ref}$ and $w_{\rm ref}$ are measured against the same scale $\sqrt{\varepsilon}$.

\subsection{Impact on the diffusion limit}

The uncoupled operational-time lattice update converges to a reaction--diffusion or fractional reaction--diffusion equation under the fixed-grid scaling
\begin{equation}
D_{\alpha_u}=\lim_{\Delta x\to0,\,\Delta u\to0}
\frac{r}{2}\frac{\Delta x^2}{(\Delta u)^{\alpha_u}}.
\end{equation}
For the coupled equation to have a finite continuum limit, the source term must remain bounded in the topology used for convergence. On a bounded local spread neighbourhood, the regularised coupling satisfies this requirement because
\begin{equation}
|\ell_{\rm reg}(y,z)|
\leq \gamma_{jk}|z|\,|q_j(y)|,
\end{equation}
and $q_j$ is integrable with finite first moment.  Therefore the coupling enters the continuum equation as an ordinary source term of order $O(z)$.

The bounded regularised source is therefore compatible with this continuum scaling. Any separate use of the inverse-spread branches in the finite-grid lattice equations would require additional uniform control near $z=0$; those branches are not used in the analytical limit.

The limiting order used in the analytical construction is therefore:
\begin{enumerate}
\item introduce the bounded source used for the continuum analytical calculation;
\item take the operational-time diffusion or grid limit with $D_{\alpha_u}$ fixed and with source moments uniformly bounded;
\item perform the local small-spread expansion of the continuum source to obtain the finite ordered-pair response coefficient $\kappa_{jk}$;
\item if desired, take a hard-side limit of the selector after the first moment has been matched.
\end{enumerate}

Conditional on the printed-equation lattice first-moment criterion, this weak local correspondence is sufficient for the log-price reduction because the Epps-effect calculation depends on the finite ordered-pair rates $\kappa_{12}$ and $\kappa_{21}$, and hence on the spread-relaxation rate $\kappa=\kappa_{12}+\kappa_{21}$. It does not require pointwise equality of the complete source fields.

\section{First-moment correspondence}
\label{app:first-moment-correspondence}

Here we distinguish the stated lattice equations in \citet{BauerDianaGebbie2024} from the actual numerical implementation\footnote{\url{https://github.com/DominicGBauer/InteractingLOBs.jl}}, and then the continuum sources we use. We prove the exact continuum first-moment identity and state a separate lattice criterion for their stated source equations \cite{BauerDianaGebbie2024}. 

\subsection{Finite-grid equations}

\citet{BauerDianaGebbie2024} write
\begin{equation}
    \Delta p_{jk}=p^{(j)}-p^{(k)}=z_{jk}
\end{equation}
and use the finite-domain source shape
\begin{equation}
    g_j(x)=-\lambda_j\mu_jx e^{-\mu_jx^2}
\end{equation}
and then state a piecewise lattice rule, denoted here by $G_{\rm B}^{(j,k)}$ to keep both ordered-pair indices visible:
{\scriptsize
\begin{equation}
G_{\rm B}^{(j,k)}=
\left\{
\begin{array}{ll}
g_j\!\left(x-p^{(j)}(t)\right)\Delta p_{jk},
&x>p^{(j)}(t),\ \Delta p_{jk}>0,\\
g_j\!\left(\dfrac{j}{\Delta p_{jk}}(x-p^{(j)}(t))\right),
&x\leq p^{(j)}(t),\ \Delta p_{jk}>0,\\
g_j\!\left(x-p^{(j)}(t)\right)\Delta p_{jk},
&x\leq p^{(j)}(t),\ \Delta p_{jk}\leq0,\\
g_j\!\left(\dfrac{1}{\Delta p_{jk}}(x-p^{(j)}(t))\right),
&x>p^{(j)}(t),\ \Delta p_{jk}\leq0.
\end{array}
\right.
\label{eq:bauer-grid-source}
\end{equation}
}
However, the public Julia implementation uses the same one-front price architecture as our analytical model: each book's mid-price is extracted from the zero crossing of the signed density, with linear interpolation between adjacent grid points. Its pair-trader coupling deforms opposite sides of that same front through $f(y/g)$ and $g f(y)$, with $g=1+c\tanh(|z|/a)$, and uses the equivalent decaying-Gaussian parameterisation $f(y)\propto-\mu y\exp[-(\mu y)^2]$. At $z=0$, $g=1$, so this coupling is a centred odd Gaussian deformation and does not by itself displace the symmetric zero crossing. We therefore do not identify its full raw first moment with the signed $O(z)$ moment below; the regularised analytical source is a reduced first-order representation of the directed front response.

\subsection{Distinct continuum sources}

Write $y=x-p^{(j)}$ and $z=z_{jk}=p^{(j)}-p^{(k)}$. The smooth continuum source is
\begin{equation}
    \ell_\varepsilon^{(j,k)}(y,z)
    =\gamma_{jk}zq_j(y)W(y,z;\varepsilon),
\end{equation}
whereas the hard-side source is
\begin{equation}
    \ell_0^{(j,k)}(y,z)
    =\gamma_{jk}zq_j(y)\mathbf 1_{\{yz>0\}}.
\end{equation}
These continuum fields are not pointwise identical to the stated symmetric lattice rule \cite{BauerDianaGebbie2024}, and for finite $\varepsilon$ the smooth and hard-side continuum fields are not pointwise identical to one another:
\begin{equation}
    G_{\rm B}^{(j,k)}\not\equiv\ell_0^{(j,k)},
    \qquad
    \ell_0^{(j,k)}\not\equiv\ell_\varepsilon^{(j,k)}.
\end{equation}

Now define the raw first-moment functional
\begin{equation}
    \mathcal M_j[\ell](z)
    =\int_{\mathbb R}y\ell(y,z)\,dy.
\end{equation}
Assume that $q_j$ is odd and integrable with finite first moment and that the local integration domain is symmetric. Then $yq_j(y)$ is even, whereas
\begin{equation}
    yq_j(y)\tanh\!\left(\frac{yz}{\varepsilon}\right)
\end{equation}
is odd in $y$. Hence, for every finite $\varepsilon>0$,
\begin{align}
\int_{\mathbb R}yq_j(y)W(y,z;\varepsilon)\,dy
&=\frac12\int_{\mathbb R}yq_j(y)\,dy\\
&=\int_0^\infty yq_j(y)\,dy
=:M_j^+.
\end{align}
The hard-side source gives the same moment for either sign of $z$. Therefore
\begin{equation}
\mathcal M_j[\ell_\varepsilon^{(j,k)}](z)
=\mathcal M_j[\ell_0^{(j,k)}](z)
=\gamma_{jk}M_j^+z.
\label{eq:exact-raw-moment}
\end{equation}
For the decaying Gaussian continuum kernel,
\begin{equation}
    q_j(y)=-\lambda_j\mu_jy e^{-\mu_jy^2},
\end{equation}
this coefficient is
\begin{equation}
    M_j^+
    =-\lambda_j\mu_j\int_0^\infty y^2e^{-\mu_jy^2}\,dy
    =-\frac{\lambda_j\sqrt\pi}{4\sqrt{\mu_j}}.
\end{equation}
Equation~\eqref{eq:exact-raw-moment} is an exact identity for the raw moment and requires no $\varepsilon\to0$ limit. Conversion of this moment into the signed front velocity uses the separately normalized closure in \ref{app:coupling-only}.

\subsection{Source-equation lattice criterion}

For lattice points $y_i=x_i-p^{(j)}$, define
\begin{equation}
    \mathcal M_{\Delta x}[G_{\rm B}](z)
    =\Delta x\sum_i y_iG_{{\rm B},i}^{(j,k)}(z).
\end{equation}
The lattice and continuum sources have a common leading raw moment if
\begin{equation}
    \mathcal M_{\Delta x}[G_{\rm B}](z)
    =\gamma_{jk}M_j^+z+r_{\Delta x}(z),
\end{equation}
with, for example, the controlled iterated condition
\begin{equation}
    \lim_{z\to0}\lim_{\Delta x\to0}
    \frac{r_{\Delta x}(z)}{z}=0,
\end{equation}
or a separately stated joint scaling. This is the matching criterion required for a common leading response coefficient after applying the same reaction-front closure. It is not a convergence theorem for the complete lattice source. The inverse-spread branches do not satisfy the bounded whole-line hypotheses used in the continuum calculation.

\subsection{Dimensions, scale and limiting order}

The selector is
\begin{equation}
    W(y,z;\varepsilon)
    =\frac12\left[1+\tanh\!\left(\frac{yz}{\varepsilon}\right)\right].
\end{equation}
Since $y$ and $z$ both have units of log-price,
\begin{equation}
    [\varepsilon]=[\mathrm{log\text{-}price}]^2.
\end{equation}
For fixed $z\neq0$, the transition width in the local log-price coordinate is
\begin{equation}
    w_\varepsilon(z)\sim\frac{\varepsilon}{|z|}.
\end{equation}
A general reference-scale choice is
\begin{equation}
    \varepsilon=|z|_{\rm ref}w_{\rm ref},
    \qquad
    \Delta x\ll w_{\rm ref}\ll L_{q,j},
    \qquad
    L_{q,j}\sim\mu_j^{-1/2}.
\end{equation}
The condition $\Delta x\ll\sqrt\varepsilon\ll L_{q,j}$ is recovered only as a symmetric reference-scale convention in which $|z|_{\rm ref}$ and $w_{\rm ref}$ are measured against the same scale $\sqrt\varepsilon$.

The pointwise limits $z\to0$ and $\varepsilon\to0$ do not commute: fixed $\varepsilon$ gives $W\to1/2$ as $z\to0$, whereas the hard-side limit at fixed nonzero $z$ gives $W\to\mathbf 1_{\{yz>0\}}$. This non-uniformity does not affect Equation~\eqref{eq:exact-raw-moment} under the stated parity and integrability assumptions. It does prevent the regularisation from being read as a pointwise rewrite of the finite-grid rule.

\section{Limiting correspondence}
\label{app:legacy-interpolation}

The analytical model presented here uses a fixed operational-time and spatial lattice and introduces calendar time by subordination. The earlier numerical papers instead evolved the densities directly on a non-uniform time grid written in the variable $t$ \cite{DianaGebbie2025} where the realised temporal increments enter the state dynamics and the pair-trader response together. This appendix records the general book-indexed construction and states the level at which it corresponds to the revised analytical model.

\subsection{Numerical construction}

For book $j$, write the non-uniform numerical grid as
\begin{equation}
    0<t_1^{(j)}<t_2^{(j)}<\cdots,
    \qquad
    \Delta t_m^{(j)}=t_{m+1}^{(j)}-t_m^{(j)}.
\end{equation}
The event-specific log-price displacement is
\begin{equation}
    \delta x_m^{(j)}
    =\left(\frac{2D_{\alpha_{\rm num}}}{r}\right)^{1/2}
    \left(\Delta t_m^{(j)}\right)^{\alpha_{\rm num}/2}.
    \label{eq:legacy-event-displacement}
\end{equation}
Here $\alpha_{\rm num}$ denotes the single fractional exponent used in the earlier numerical construction. Let $\mathcal I_{\Delta x}$ denote the interpolation operator on the background log-price lattice with spacing $\Delta x$. The hatted off-grid neighbouring states are
\begin{equation}
    \widehat\varphi_m^{\,i\pm1(j)}
    =\mathcal I_{\Delta x}\varphi_\Delta^{(j)}
    \left(x_i\pm\delta x_m^{(j)},t_m^{(j)}\right).
    \label{eq:legacy-interpolated-neighbours}
\end{equation}
Writing the numerical pair-trader source as $\ell_{{\rm num},m}^{i(j,k)}$, the source increment is
\begin{equation}
    c_m^{i(j)}=s_m^{i(j)}+\delta_m^{i(j)}
    +\sum_{k\ne j}\ell_{{\rm num},m}^{i(j,k)},
\end{equation}
and the book-indexed non-uniform update is
\begin{equation}
\begin{aligned}
\varphi_n^{i(j)}
&=\sum_{m=0}^{n-1}K_{n-m}^{(\alpha_{\rm num})}
 e^{-\nu_j(t_{n-1}^{(j)}-t_m^{(j)})}\bigg[
 \tfrac12(r+F_{n-1}^{(j)})\widehat\varphi_m^{\,i-1(j)}\\
&\quad+\tfrac12(r-F_{n-1}^{(j)})\widehat\varphi_m^{\,i+1(j)}
-r\varphi_m^{i(j)}\bigg]\\
&\quad+e^{-\nu_j\Delta t_{n-1}^{(j)}}\varphi_{n-1}^{i(j)}
+c_{n-1}^{i(j)}\Delta t_{n-1}^{(j)}.
\end{aligned}
\label{eq:legacy-interpolated-update}
\end{equation}
Equations~\eqref{eq:legacy-event-displacement}--\eqref{eq:legacy-interpolated-update} give the general form of \citet{DianaGebbie2025,BauerDianaGebbie2024}. 

\subsection{Common law and common realisation}

The books may use the same waiting-time law and average event scale,
\begin{equation}
    \Delta t_m^{(1)}\sim\Psi,
    \qquad
    \Delta t_m^{(2)}\sim\Psi,
    \qquad
    \mathbb E[\Delta t_m^{(1)}]
    =\mathbb E[\Delta t_m^{(2)}],
\end{equation}
while using different realised sequences pathwise. A common realised grid is the special case
\begin{equation}
    t_m^{(1)}=t_m^{(2)}=t_m.
\end{equation}
The analytical argument below does not depend on resolving which special case was used in the earlier code.

\subsection{Mesh and coupling entanglement}

The realised increment $\Delta t_m^{(j)}$ determines the off-grid displacement,
\begin{equation}
    \delta x_m^{(j)}
    \propto
    \left(\Delta t_m^{(j)}\right)^{\alpha_{\rm num}/2},
\end{equation}
and also enters the cancellation survival factors, source accumulation and hatted interpolated states. The pair-trader source is evaluated from order-density fields and reaction boundaries whose paths have already been changed by those increments. The numerical construction therefore does not admit a pathwise factorisation into an independent clock term and an independent coupling term.

If the books share one realised non-uniform grid, temporal propagation and coupling response remain entangled because the same random increments enter the state update. What disappears is only the additional loss caused by incomplete overlap between distinct realised observation clocks. If the books use separate realisations, the numerical construction may contain both the entangled non-uniform state evolution and cross-book clock asynchrony.

\subsection{Continuum correspondence}

The non-uniform numerical method and the present construction are not pointwise or pathwise identical. Suppose that the centred case $F_m^{(j)}=0$ is used and that the interpolated density is locally smooth. Let $\mathcal I_{\Delta x}$ be consistent and stable. For book $j$,
\begin{equation}
\begin{aligned}
&\mathcal I_{\Delta x}\varphi^{(j)}
\left(x_i\pm\delta x_m^{(j)},t_m^{(j)}\right)\\
&\quad=\varphi^{(j)}\left(x_i,t_m^{(j)}\right)
\pm\delta x_m^{(j)}\partial_x\varphi^{(j)}
\left(x_i,t_m^{(j)}\right)\\
&\qquad+\tfrac{1}{2}(\delta x_m^{(j)})^2
\partial_x^2\varphi^{(j)}\left(x_i,t_m^{(j)}\right)
+R_{i,m}^{\pm(j)}.
\end{aligned}
\end{equation}
The symmetric nearest-neighbour combination cancels the first derivative and gives
\begin{equation}
\begin{aligned}
&\frac{r}{2}\left(
\widehat\varphi_m^{\,i-1(j)}
+\widehat\varphi_m^{\,i+1(j)}
-2\varphi_m^{i(j)}\right)\\
&\quad=\frac{r}{2}(\delta x_m^{(j)})^2
\partial_x^2\varphi^{(j)}\left(x_i,t_m^{(j)}\right)
+R_{i,m}^{(j)}.
\end{aligned}
\end{equation}
Using Equation~\eqref{eq:legacy-event-displacement},
\begin{equation}
    \frac{r}{2}\frac{(\delta x_m^{(j)})^2}
    {(\Delta t_m^{(j)})^{\alpha_{\rm num}}}
    =D_{\alpha_{\rm num}}.
\end{equation}
A sufficient correspondence condition is that the accumulated interpolation and Taylor remainders vanish in the weak limit. One convenient local condition is
\begin{equation}
    \frac{|R_{i,m}^{(j)}|}
    {(\Delta t_m^{(j)})^{\alpha_{\rm num}}}
    \longrightarrow0
\end{equation}
uniformly on the compact space--time region used in the test-function argument. For piecewise-linear interpolation, a stronger but transparent sufficient scaling is that the background-grid interpolation error be of smaller order than $(\delta x_m^{(j)})^2$. More general joint grid and waiting-time limits may be used if the weighted remainder is controlled directly.

Under these assumptions, the leading centred transport term approximates a fractional transport law in the original numerical time variable $t$. Its correspondence with the present operational-time model is obtained only after the operational law is projected into calendar time by subordination and the resulting continuum or weak equations are matched.

\subsection{Relating the clocks}

The prior numerical simulation exponent $\alpha_{\rm num}$ from \cite{BauerDianaGebbie2024} is an effective exponent and cannot be identified uniquely with either $\alpha_u$ or $\alpha_t$ without an additional modelling restriction. With an identity observation clock, $E_j(t)=t$, it may be identified with the operational-time exponent, $\alpha_{\rm num}=\alpha_u$. If the operational dynamics is ordinary, $\alpha_u=1$, and anomalous scaling is assigned entirely to the inverse clock, then $\alpha_{\rm num}=\alpha_t$. When both mechanisms are retained, no general one-to-one identification is imposed.

The analytical treatment therefore isolates a clock-overlap factor and a finite coupling-response factor from mechanisms that are entangled in the numerical state update -- this is a key contribution in this work. A common analytical observation clock, $E_1=E_2$, removes only the separate clock-overlap attenuation; it does not imply that a non-uniform numerical state evolution was free of temporal effects.


\section{Analytical clock and coupling-response separation}
\label{app:analytical-numerical-correspondence}

The earlier numerical model generated the Epps effect from a coupled, non-uniform event-time construction in which the clock and coupling effects were entangled. The present analytical treatment separates these contributions into a clock-overlap factor and a finite coupling-response factor. Their leading-order product provides a reduced description of the combined mechanism; it is not intended as a reproduction or fit of the published numerical curve \citep{BauerDianaGebbie2024}.

Figure~\ref{fig:analytical-numerical-correspondence} visualises the separated clock, coupling-response and combined factors using effective rates:
\begin{equation}
\frac{\rho_{\Delta t}^{\mathrm{comb}}}{\rho_\infty}
\simeq
\mathcal A_{12}(\Delta t)\mathcal R_\kappa(\Delta t),
\qquad
\mathcal R_\kappa(\Delta t)
=1-\frac{1-e^{-\kappa\Delta t}}{\kappa\Delta t},
\label{eq:analytical-numerical-correspondence}
\end{equation}
where the exponential clock-overlap envelope used is:
\begin{equation}
\mathcal A_{12}(\Delta t)
=1-\frac{1-e^{-\lambda_{12}^{\rm clk}\Delta t}}
{\lambda_{12}^{\rm clk}\Delta t}.
\end{equation}
Here $\kappa$ is the effective calendar-time coupling-response rate for the exponential or identity-clock benchmark. Under non-trivial inverse-stable subordination, the corresponding response is Mittag--Leffler, as in \ref{app:combined-subordination-coupling}.

\begin{figure}[htbp]
\centering
\includegraphics[width=\columnwidth]{figures/figure-06-calendar-time-epps-memory-v1.pdf}
\caption{\textbf{Analytical separation of clock and coupling-response effects.} The clock-overlap factor (black), coupling-response factor (blue) and their leading-order separable product (red) are shown using the illustrative rates $\lambda_{12}^{\rm clk}=0.1\,\mathrm{s}^{-1}$ and $\kappa=0.025\,\mathrm{s}^{-1}$, with $\rho_\infty=1$. The main panel is normalised by the limiting correlation. The inset shows the corresponding no-refresh and coupling-relaxation survival functions. The main horizontal axis gives the fixed calendar-time aggregation interval $\Delta t$, not the indexed numerical increment $\Delta t_m^{(j)}$. The inset uses operational lag $u$. The figure contextualises the analytical decomposition; it is not a reproduction, fit or parameter-by-parameter validation of the published numerical result.}
\label{fig:analytical-numerical-correspondence}
\end{figure}

\end{document}
